\documentclass[11pt]{article}
\usepackage[margin=1.1in]{geometry}
\usepackage{amsmath,amssymb,amsthm}
\usepackage[colorlinks=true,linkcolor=blue,citecolor=blue,urlcolor=blue]{hyperref}
\usepackage{enumitem}

\newtheorem{theorem}{Theorem}
\newtheorem{proposition}{Proposition}
\newtheorem{lemma}{Lemma}
\newtheorem{corollary}{Corollary}
\newtheorem{assumption}{Assumption}
\newtheorem{conjecture}{Conjecture}
\theoremstyle{remark}
\newtheorem{remark}{Remark}

\newcommand{\Cam}{C^{\mathrm{am}}}
\newcommand{\Ceu}{C^{\mathrm{eu}}}
\newcommand{\E}{\mathbb{E}}
\newcommand{\Prob}{\mathbb{P}}
\newcommand{\Kt}[1]{\tilde K_{#1}}

\title{TP\textsubscript{2} for American Calls:\\ Reductions, a Dividend-Robust Violation Certificate,\\ and a Supermodular Optimal-Stopping Route\thanks{Companion theory note to the ICAIF~'26 submission \emph{Machine-Learned Conditional Trading of TP2 Violations in SPY American Call Options}. Propositions~\ref{prop:window} and~\ref{prop:certificate} and Lemma~\ref{lem:smalldiv} are stated with complete proofs for latent values; Conjecture~\ref{thm:topkis} is a proof \emph{program} with the remaining gaps stated explicitly.}}
\author{}
\date{August 2, 2026}

\begin{document}
\maketitle

\begin{abstract}
\noindent The TP\textsubscript{2} property of call prices---$C(K_1,T_1)C(K_2,T_2)\ge C(K_1,T_2)C(K_2,T_1)$ for $K_1\le K_2$, $T_1\le T_2$---is proved by Glasserman and Pirjol for \emph{European} options. SPY calls are American-style on a dividend-paying ETF, so the theorem does not directly apply. This note shows the gap is smaller than it looks. (i)~If no ex-dividend date falls before expiry, or if a simple rate-dependent condition on dividend sizes holds, then American SPY calls \emph{equal} their European counterparts and the TP\textsubscript{2} theory applies as is. Higher interest rates expand the set of dividend-spanning contracts satisfying the condition; whether it holds is contract-specific, and at near-zero rates it fails for contracts spanning an ex-date. (ii)~When early exercise cannot be ruled out, the early-exercise premium is at most the present value of the interim dividends. This gives a computable margin: any executable-side TP\textsubscript{2} violation larger than the margin implies a violation of the \emph{European} inequality for latent values. The margin has the units of the price products; normalized versions ($M/S_0^2$, or $M$ relative to the overvalued product) are the economically interpretable quantities, with each dividend term of order $0.3\%$--$0.6\%$ of the underlying price. (iii)~Concentration bounds sharpen the margin. The premium is weighted by the probability of exercise. Doob's maximal inequality bounds that probability from quotes and martingality alone; a separate Chernoff bound requires a sub-Gaussian modeling assumption, with an implied-volatility proxy as one calibration. They cut the margin by an order of magnitude for deep out-of-the-money, short-dated rectangles, which is where violations concentrate. A perturbation lemma then treats the American premium, bid/ask spreads, and strike rounding as one error term. (iv)~We outline a route to American TP\textsubscript{2} without any reduction, via supermodularity of the optimal-stopping value (Topkis). In the discrete-dividend case this reduces to finitely many fixed-date comparisons and scale maps that preserve the Glasserman--Pirjol tail condition. We state which steps remain open.
\end{abstract}

\section{Setting and Notation}
Fix a current date $t=0$, a filtered probability space carrying the (risk-neutral) dynamics of the underlying $S=\{S_u\}_{u\ge0}$, and a constant rate $r\ge0$. Dividends are discrete: ex-dividend dates $0<t_1<t_2<\cdots$ with per-share amounts $\delta_1,\delta_2,\ldots\ge 0$ (SPY pays quarterly; near-term amounts are announced in advance and we treat them as known; unannounced far dividends can be handled by replacing $\delta_i$ with a conservative upper estimate wherever an upper bound is used, and a lower estimate wherever a lower bound is used). For strike $K>0$ and expiry $T>0$ write $\Ceu(K,T)$ and $\Cam(K,T)$ for European and American call prices, and
\[
\bar D(T)\;=\;\sum_{i:\,0<t_i\le T}\delta_i\,e^{-r t_i},
\qquad
\pi(K,T)\;=\;\Cam(K,T)-\Ceu(K,T)\;\ge0 .
\]
The forward-adjusted TP\textsubscript{2} condition of Glasserman--Li--Pirjol, for $T_1<T_2$ and $K_1/F_{T_1}<K_2/F_{T_2}$ with forwards $F_T$ and adjusted strikes $\Kt1=K_1F_{T_2}/F_{T_1}$, $\Kt2=K_2F_{T_1}/F_{T_2}$ (rounded up to traded strikes), reads
\begin{equation}
C(K_1,T_1)\,C(K_2,T_2)\;\ge\;C(\Kt1,T_2)\,C(\Kt2,T_1).
\label{eq:tp2adj}
\end{equation}
Glasserman--Pirjol (\emph{Math.\ Finance}, 2026) prove \eqref{eq:tp2adj} for European prices under distributional conditions on $(S_{T_1},S_{T_2})$: a conditional right-tail dominance condition (TP\textsubscript{2} of the survival function) on the out-of-the-money region, and an increasing-convex-order condition with a left-tail comparison covering in-the-money strikes (their Theorem~3.1). We refer to these jointly as \emph{the GP conditions}. Throughout, ``the four legs'' of a rectangle means the contracts at $(K_1,T_1),(K_2,T_2),(\Kt1,T_2),(\Kt2,T_1)$.

Two classical facts, both due to Merton~\cite{merton}, with the second stated in the exact form we use:
\begin{lemma}[Merton reduction~\cite{merton}]\label{lem:merton}
If no dividend is paid in $(0,T]$, then $\Cam(K,T)=\Ceu(K,T)$ for every $K$.
\end{lemma}
\begin{lemma}[Dividend sandwich]\label{lem:sandwich}
In general, $\;0\le \pi(K,T)\le \bar D(T)$, i.e.\ $\Ceu(K,T)\le\Cam(K,T)\le \Ceu(K,T)+\bar D(T)$.
\end{lemma}
\begin{proof}[Proof of Lemma~\ref{lem:sandwich}]
The lower bound is immediate (the American holder can always wait to $T$). For the upper bound, hold the European call plus a claim paying each dividend $\delta_i$ at $t_i$ (cost $\Ceu+\bar D(T)$), against a short American call. If the American holder exercises at some $\tau$ with $S_\tau-K>0$, the long portfolio is worth $\Ceu(S_\tau,\tau)+\sum_{t_i>\tau}\delta_i e^{-r(t_i-\tau)}$ plus the accrued dividends, which dominates $S_\tau-K$ by the European lower bound on the dividend-adjusted forward; if never exercised, both calls pay $(S_T-K)^+$. Either way the long side is never worse, so $\Cam\le\Ceu+\bar D(T)$; the argument is Merton's~\cite{merton}.
\end{proof}

\section{Europeanization: When American SPY Calls Are European}

\begin{lemma}[Rate-dependent small-dividend Europeanization]\label{lem:smalldiv}
Suppose that for every ex-dividend date $t_i\in(0,T]$,
\begin{equation}
\delta_i\;+\;\sum_{j:\,t_i<t_j\le T}\delta_j\,e^{-r(t_j-t_i)}\;\le\;K\left(1-e^{-r(T-t_i)}\right).
\label{eq:smalldiv}
\end{equation}
Then early exercise of the American call $(K,T)$ is never optimal and $\Cam(K,T)=\Ceu(K,T)$.
\end{lemma}

\begin{proof}
Early exercise of an American call on a discrete-dividend asset can be optimal only at an instant immediately preceding an ex-dividend date (between dividends the argument of Lemma~\ref{lem:merton} applies on the dividend-free segment). Consider exercise at $t_i^-$, when the cum-dividend price is $S$. Exercising yields $S-K$. Not exercising leaves a call on the ex-dividend asset worth, by the forward lower bound with discrete dividends (the same no-arbitrage portfolio argument as Lemma~\ref{lem:sandwich}),
\[
\Cam\big(S-\delta_i,\,t_i;K,T\big)\;\ge\;(S-\delta_i)\;-\;\sum_{j>i,\,t_j\le T}\delta_j e^{-r(t_j-t_i)}\;-\;Ke^{-r(T-t_i)} .
\]
Continuation dominates exercise whenever the right-hand side is at least $S-K$, which rearranges to \eqref{eq:smalldiv}. If \eqref{eq:smalldiv} holds at every $t_i$, no exercise date is ever strictly preferred, so the American and European values coincide. For a single interim dividend, \eqref{eq:smalldiv} reduces to the classical criterion $\delta\le K(1-e^{-r(T-\tau)})$ underlying Black's approximation~\cite{black1975}.
\end{proof}

\begin{remark}[The regime content of \eqref{eq:smalldiv}]\label{rem:regime}
For $K\approx S$, a dividend of approximately $0.31\%K$, and 60 calendar days from the ex-date to expiry, the right-hand side of \eqref{eq:smalldiv} is approximately $0.82\%K$ at $r=5\%$ ($1-e^{-0.05\cdot 60/365}\approx0.8185\%$) and $0.16\%K$ at $r=1\%$ ($1-e^{-0.01\cdot 60/365}\approx0.1642\%$). The condition is therefore satisfied in the first case but not the second. The break-even horizon for a single dividend is
\[
T-t_i\;\ge\;-\frac{\log(1-\delta_i/K)}{r},
\]
approximately $55$ calendar days at $r=2\%$ and $22$ calendar days at $r=5\%$ for $\delta_i/K=0.003$. Higher interest rates expand the set of dividend-spanning contracts satisfying the sufficient Europeanization condition. Whether the condition holds remains contract-specific and depends on the strike, dividend amount, and time from each ex-date to expiry. At $r\approx0$ (the 2021 regime) the right-hand side vanishes and the condition fails for every contract spanning an ex-date. Deep in-the-money strikes (small $K$) tighten \eqref{eq:smalldiv}, consistent with early exercise being an in-the-money phenomenon.
\end{remark}

\begin{proposition}[No-dividend-window / certified-window TP\textsubscript{2}]\label{prop:window}
Fix a violation rectangle with far maturity $T_2$. If either (a) $(0,T_2]$ contains no ex-dividend date, or (b) condition \eqref{eq:smalldiv} holds for each of the four legs, then the latent American no-arbitrage value of each contract equals its European value, and under the GP conditions on $(S_{T_1},S_{T_2})$ these latent values satisfy the forward-adjusted TP\textsubscript{2} inequality \eqref{eq:tp2adj}. Consequently, an executable-side violation on such a rectangle (asks on the undervalued pair below bids on the overvalued pair, so that the latent values are bracketed against the inequality) contradicts the \emph{European} TP\textsubscript{2} theory itself, exactly as on SPX.
\end{proposition}

\begin{proof}
Immediate from Lemma~\ref{lem:merton} (case a) or Lemma~\ref{lem:smalldiv} (case b) applied to each leg, followed by Glasserman--Pirjol Theorem~3.1 in its forward-adjusted form.
\end{proof}

\section{A Dividend-Robust Violation Certificate}

When neither reduction applies, the American premium $\pi\ge0$ perturbs \eqref{eq:tp2adj} in a direction that is \emph{not} sign-definite (Section~\ref{sec:direction}), so an observed violation of American quotes does not by itself contradict the European theory. It does, however, whenever the violation exceeds an explicit margin.

\begin{proposition}[Certificate]\label{prop:certificate}
Let $\bar D_1=\bar D(T_1)$ and $\bar D_2=\bar D(T_2)$. If quoted American prices satisfy
\begin{equation}
\Cam(K_1,T_1)\,\Cam(K_2,T_2)\;<\;\Cam(\Kt1,T_2)\,\Cam(\Kt2,T_1)\;-\;M,
\qquad
M:=\bar D_2\,\Cam(\Kt2,T_1)+\bar D_1\,\Cam(\Kt1,T_2),
\label{eq:certificate}
\end{equation}
then the corresponding European prices violate \eqref{eq:tp2adj}: $\Ceu(K_1,T_1)\Ceu(K_2,T_2)<\Ceu(\Kt1,T_2)\Ceu(\Kt2,T_1)$. In particular, under the GP conditions such a configuration is impossible in any model in which they hold, so a market quote satisfying \eqref{eq:certificate} is a genuine anomaly not attributable to early-exercise premia.
\end{proposition}

\begin{proof}
By Lemma~\ref{lem:sandwich}, $\Ceu\le\Cam$ pointwise, so the left side satisfies
$\Ceu(K_1,T_1)\Ceu(K_2,T_2)\le \Cam(K_1,T_1)\Cam(K_2,T_2)$.
Also by Lemma~\ref{lem:sandwich}, $\Ceu\ge\Cam-\bar D(T)$ at each leg, hence
\begin{align*}
\Ceu(\Kt1,T_2)\,\Ceu(\Kt2,T_1)
&\ge \big(\Cam(\Kt1,T_2)-\bar D_2\big)\big(\Cam(\Kt2,T_1)-\bar D_1\big)\\
&= \Cam(\Kt1,T_2)\Cam(\Kt2,T_1)-\bar D_2\Cam(\Kt2,T_1)-\bar D_1\Cam(\Kt1,T_2)+\bar D_1\bar D_2\\
&\ge \Cam(\Kt1,T_2)\Cam(\Kt2,T_1)-M .
\end{align*}
Chaining the two displays with \eqref{eq:certificate} gives
$\Ceu(K_1,T_1)\Ceu(K_2,T_2)\le \Cam(K_1,T_1)\Cam(K_2,T_2)<\Cam(\Kt1,T_2)\Cam(\Kt2,T_1)-M\le \Ceu(\Kt1,T_2)\Ceu(\Kt2,T_1)$.
\end{proof}

\begin{remark}[Units and normalization of the margin]\label{rem:units}
For the maturities in the sample, each $\bar D_i$ is typically of the order of one or two quarterly SPY dividends, approximately $0.3\%$--$0.6\%$ of the contemporaneous underlying price, and $\bar D_1=0$ whenever no ex-date precedes $T_1$. The certificate margin $M$, however, has the same units as the TP\textsubscript{2} price products (currency squared). Its economic magnitude should therefore be reported after normalization, for example
\[
m_S=\frac{M}{S_0^2},\qquad
m_R=\frac{M}{\Cam(\Kt1,T_2)\,\Cam(\Kt2,T_1)},
\]
or as the equivalent log-defect threshold ($m_R$ is exactly the relative margin that the normalized violation premium must exceed). $M$ is computable trade by trade from announced dividends, and if leg $\ell$ satisfies Lemma~\ref{lem:smalldiv}, its $\bar D$ can be replaced by $0$.
\end{remark}

\begin{remark}[Latent values versus quotes; the executable certificate]\label{rem:quotes}
Lemmas~\ref{lem:merton}--\ref{lem:smalldiv} and Proposition~\ref{prop:certificate} are statements about latent no-arbitrage values. Write $C^{\mathrm{am},*}$ for the latent American value and $C^{\mathrm{bid}}\le C^{\mathrm{am},*}\le C^{\mathrm{ask}}$ for the standing quotes (an assumption on quote quality, not a theorem). The executable version of the certificate then uses directional quote bounds: on an undervalued-side contract, $\Ceu\le C^{\mathrm{am},*}\le C^{\mathrm{ask}}$; on an overvalued-side contract, $\Ceu\ge C^{\mathrm{am},*}-\bar D\ge C^{\mathrm{bid}}-\bar D$. Consequently, if
\[
C^{\mathrm{ask}}(K_1,T_1)\,C^{\mathrm{ask}}(K_2,T_2)\;<\;C^{\mathrm{bid}}(\Kt1,T_2)\,C^{\mathrm{bid}}(\Kt2,T_1)\;-\;M^{\mathrm{bid}},
\qquad
M^{\mathrm{bid}}=\bar D_2\,C^{\mathrm{bid}}(\Kt2,T_1)+\bar D_1\,C^{\mathrm{bid}}(\Kt1,T_2),
\]
then the latent European values violate \eqref{eq:tp2adj}, by the same three-inequality chain. This matches the empirical pipeline's detection convention (undervalued product at asks, overvalued at bids). We call violations satisfying this executable inequality \emph{dividend-robust}.
\end{remark}

\section{Sharper Margins: Concentration and Perturbation Stability}\label{sec:concentration}

The certificate of Proposition~\ref{prop:certificate} uses the worst-case bound $\pi\le\bar D$, which charges the full present value of dividends to every contract regardless of strike. This is loose in exactly the region where violations concentrate: the SPX evidence locates most violations at fairly deep out-of-the-money strikes and short expiries, where the probability of the underlying ever reaching the exercise region is small. This section replaces the worst case by exercise-probability-weighted bounds---first exactly, then via two computable relaxations (one model-free from quoted prices via Doob's maximal inequality, one from an implied-volatility proxy via a Chernoff bound)---and closes with a perturbation lemma showing that the TP\textsubscript{2} defect, viewed in log form, is $1$-Lipschitz in \emph{every} friction at once.

\subsection{The Premium Is an Exercise-Probability-Weighted Quantity}

\begin{proposition}[Exercise-probability premium bound]\label{prop:probweight}
Suppose a single ex-dividend date $t_D\in(0,T]$ with amount $\delta$. Then
\begin{equation}
\pi(K,T)\;\le\;\delta\,e^{-rt_D}\,\Prob\!\left(S_{t_D^-}\ge K\right).
\label{eq:probweight}
\end{equation}
With $n$ interim dividends, $\pi(K,T)\le\sum_i w_i\,e^{-rt_i}\,\Prob(S_{t_i^-}\ge K)$ with weights $w_i=\delta_i+\sum_{j>i,\,t_j\le T}\delta_j e^{-r(t_j-t_i)}$.
\end{proposition}

\begin{proof}
Single dividend: exercise can be optimal only at $t_D^-$, and the post-dividend continuation is European (no dividends remain), so the dynamic-programming identity is exact:
\[
\Cam(K,T)=\E\!\left[e^{-rt_D}\max\!\big\{S_{t_D^-}-K,\;\Ceu(S_{t_D^-}-\delta,t_D;K,T)\big\}\right],\quad
\Ceu(K,T)=\E\!\left[e^{-rt_D}\,\Ceu(S_{t_D^-}-\delta,t_D;K,T)\right].
\]
Hence $\pi=\E[e^{-rt_D}(S_{t_D^-}-K-\Ceu(S_{t_D^-}-\delta,t_D))^+]$. Pathwise, write $S=S_{t_D^-}$: if $S<K$ the integrand is $0$ (since $\Ceu\ge0$). If $S\ge K$, then using $\Ceu(S-\delta,t_D)\ge\max\{0,(S-\delta)-Ke^{-r(T-t_D)}\}$: on the branch where the max is the forward bound, $S-K-\Ceu\le \delta-K(1-e^{-r(T-t_D)})\le\delta$; on the branch where the max is $0$ we have $S\le \delta+Ke^{-r(T-t_D)}$, so $S-K\le\delta-K(1-e^{-r(T-t_D)})\le\delta$. In both cases the integrand is at most $\delta\,\mathbf 1\{S\ge K\}$, and \eqref{eq:probweight} follows. The multi-dividend statement follows by backward induction on ex-dates: at $t_i^-$, $\max\{S-K,V^{\mathrm{am}}\}-V^{\mathrm{eu}}\le (S-K-V^{\mathrm{eu}})^++(V^{\mathrm{am}}-V^{\mathrm{eu}})$, the first term is bounded pathwise by $w_i\mathbf 1\{S\ge K\}$ exactly as above (the European continuation lower bound now subtracts the remaining dividends, which is where $w_i$ acquires its second summand), and the second term recurses.
\end{proof}

Inequality \eqref{eq:probweight} is already strict progress: for an out-of-the-money leg, $\Prob(S_{t_i^-}\ge K)$ is the risk-neutral probability of moving through the strike \emph{before the ex-date}---typically far below one at the short horizons where violations cluster. It remains to bound this probability from observables.

\subsection{A Model-Free Bound from Quotes (Doob)}

\begin{proposition}[Quote-implied concentration bound, $r=0$]\label{prop:doob}
Let $r=0$ and let $X_u=S_u+\sum_{t_j\le u}\delta_j$ (the dividend-reinvested price), a $\Prob$-martingale. Fix an ex-date $t_i$, a listed maturity $\tau\ge t_i$, and any $\lambda$ with $\sum_{t_j\le\tau}\delta_j<\lambda<K_i:=K+\sum_{t_j<t_i}\delta_j$ (the lower restriction keeps the shifted strike positive). Then
\begin{equation}
\Prob\!\left(S_{t_i^-}\ge K\right)\;\le\;\Prob\!\Big(\sup_{u\le t_i}X_u\ge K_i\Big)\;\le\;\frac{\E[(X_\tau-\lambda)^+]}{K_i-\lambda}\;=\;\frac{C^{\mathrm{eu}}\!\big(\lambda-{\textstyle\sum_{t_j\le\tau}}\delta_j,\;\tau\big)}{K_i-\lambda}\;\le\;\frac{C^{\mathrm{ask}}\!\big(\lambda-{\textstyle\sum_{t_j\le\tau}}\delta_j,\;\tau\big)}{K_i-\lambda},
\label{eq:doob}
\end{equation}
optimizable over admissible $\lambda$ and over listed $\tau\ge t_i$. The final step uses $\Ceu\le C^{\mathrm{am},*}\le C^{\mathrm{ask}}$; when the shifted strike is not listed, round it \emph{down} to the nearest listed strike, which is conservative because the call price is nonincreasing in strike. The positive-rate version is derived in Appendix~\ref{app:doobr}.
\end{proposition}

\begin{proof}
$\{S_{t_i^-}\ge K\}\subseteq\{X_{t_i^-}\ge K_i\}\subseteq\{\sup_{u\le t_i}X_u\ge K_i\}$. For fixed $\lambda$, $(X_u-\lambda)^+$ is a nonnegative submartingale (convex function of a martingale), so Doob's maximal inequality gives $\Prob(\sup_{u\le\tau}(X_u-\lambda)^+\ge K_i-\lambda)\le\E[(X_\tau-\lambda)^+]/(K_i-\lambda)$, and $\sup_{u\le t_i}\le\sup_{u\le\tau}$. Finally $(X_\tau-\lambda)^+=(S_\tau-(\lambda-\sum_{t_j\le\tau}\delta_j))^+$, whose expectation is the undiscounted European call at the (positive) shifted strike, bounded above by the ask quote.
\end{proof}

Everything on the right of \eqref{eq:doob} is a market quote: the bound's ``variance proxy'' is the option surface itself.

\subsection{A Sub-Gaussian Bound from Implied Volatility}

\begin{proposition}[Model-dependent Chernoff bound]\label{prop:chernoff}
Assume explicitly that the risk-neutral log-forward return $\log(S_t/F_t)$ is sub-Gaussian with variance proxy $v_t$ (in the sense of~\cite{vershynin}, \S2.6). Then for $K\ge F_{t_i}$,
\begin{equation}
\Prob(S_{t_i}\ge K)\;\le\;\exp\!\left(-\frac{\log^2(K/F_{t_i})}{2\,v_{t_i}}\right).
\label{eq:chernoff}
\end{equation}
Time is measured in calendar years throughout this bound.
\end{proposition}
\noindent A convenient modeling choice is $v_t=\bar\sigma^2 t$ with $\bar\sigma$ inferred from the observed implied-volatility surface. This is an additional modeling approximation rather than a model-free consequence of option quotes: an implied-volatility bound is not by itself a moment-generating-function bound for the risk-neutral distribution. Numerically, using calendar-year time, a $12\%$-OTM strike at a 14-calendar-day horizon with $\bar\sigma=20\%$ gives a bound of approximately $e^{-4.18}\approx0.015$, while a $5\%$-OTM strike at a 30-calendar-day horizon gives approximately $0.70$. The first case provides substantial sharpening; the second provides only moderate sharpening.

\begin{corollary}[Two sharpened margins]\label{cor:sharpcert}
Proposition~\ref{prop:certificate} holds with $M$ replaced by either of
\[
M^{\mathrm{quote}}\;=\;\bar\pi^{\mathrm{quote}}(\Kt1,T_2)\,C^{\mathrm{bid}}(\Kt2,T_1)\;+\;\bar\pi^{\mathrm{quote}}(\Kt2,T_1)\,C^{\mathrm{bid}}(\Kt1,T_2),
\qquad
\bar\pi^{\mathrm{quote}}(K,T)\;=\;\sum_{t_i\le T} w_i\,e^{-rt_i}\,q_i^{\mathrm{quote}}(K),
\]
with $q_i^{\mathrm{quote}}(K)=\min\{1,\ \text{bound \eqref{eq:doob}}\}$, or the analogous $M^{\mathrm{model}}$ built from $q_i^{\mathrm{model}}(K)=\min\{1,\ \text{bound \eqref{eq:doob}},\ \text{bound \eqref{eq:chernoff}}\}$; $w_i$ are the weights of Proposition~\ref{prop:probweight}. $M^{\mathrm{quote}}$ uses only quotes and martingality and inherits the certificate's rigor. $M^{\mathrm{model}}$ additionally uses the sub-Gaussian assumption of Proposition~\ref{prop:chernoff}; a violation exceeding $M^{\mathrm{model}}$ but not $M^{\mathrm{quote}}$ should be labeled model-assisted or diagnostic, not certified. Since $q_i\le1$, both margins are at most $M$, with the largest improvement for deep out-of-the-money, short-dated rectangles.
\end{corollary}

\subsection{Portfolio-Level Europeanization with High Probability}

The bounds above are per contract; the trading application holds a \emph{book} of short legs. A union bound over $N$ positions would cost a factor $N$ (or $\sqrt{\log N}$ in expectation, via the $\psi_2$ maximal inequality of~\cite{vershynin}, Eq.~(2.22))---but the common-driver structure of the problem collapses it.

\begin{proposition}[Common-driver bound]\label{prop:portfolio}
Consider any set of selected trades whose short legs, alive at ex-dividend date $t_D$, have exercise regions $E_\ell\subseteq[0,\infty)$ and strikes $\{K_\ell\}$, and note that rational exercise of leg $\ell$ requires $S_{t_D^-}\ge K_\ell$ (exercise below the strike is dominated). At a fixed ex-dividend date, the event that at least one position is inside its rational exercise region can be bounded through the minimum exercise boundary across the positions:
\[
\Prob\big(\text{any short leg rationally exercisable at } t_D\big)\;\le\;\Prob\!\left(S_{t_D^-}\ge \min_\ell K_\ell\right),
\]
one application of \eqref{eq:doob} or \eqref{eq:chernoff} at the book's minimum short strike, with no union-bound inflation because all legs share one underlying. (The event equals $\{S_{t_D^-}\ge\min_\ell b_\ell\}$ for exercise boundaries $b_\ell$ only if every exercise region is a threshold region of the form $[b_\ell,\infty)$, which we do not prove here; the displayed inequality does not require it.) On the complement, no rational exercise occurs at that date: American and European portfolio cash flows coincide \emph{through} $t_D$, but future early-exercise optionality may remain. Across the (at most two or three) distinct ex-dates in a holding window, a union bound over dates costs only that small factor.
\end{proposition}

\begin{remark}[Rational exercise versus operational assignment]
Actual assignment can differ from the model-implied rational exercise rule and may alter stock exposure, financing, dividend obligations, and transaction costs. The probability bounds above therefore address rational early exercise, not the complete operational assignment process.
\end{remark}

\subsection{Perturbation Stability: One Modulus for Every Friction}

\begin{lemma}[Log-Lipschitz stability of the TP\textsubscript{2} defect]\label{lem:stability}
Define the log-defect of a rectangle as $\Delta:=\log C(K_1,T_1)+\log C(K_2,T_2)-\log C(\Kt1,T_2)-\log C(\Kt2,T_1)$ (so TP\textsubscript{2} is $\Delta\ge0$). If quoted prices satisfy $\tilde C_\ell=C_\ell e^{\eta_\ell}$ with $|\eta_\ell|\le\varepsilon_\ell$ on each leg $\ell$, then
\[
\tilde\Delta\;\ge\;\Delta\;-\;\sum_{\ell=1}^4\varepsilon_\ell .
\]
In particular, a quoted defect $\tilde\Delta<-\sum_\ell\varepsilon_\ell$ certifies $\Delta<0$ for the reference prices.
\end{lemma}
\begin{proof}
Immediate: $\tilde\Delta=\Delta+\eta_1+\eta_2-\eta_3-\eta_4\ge\Delta-\sum|\eta_\ell|$.
\end{proof}

\begin{remark}[Instantiations]\label{rem:instantiations}
Lemma~\ref{lem:stability} is the ``approximation theory'' of the problem: log-supermodularity has Lipschitz constant $1$ in each leg's log-price, so every friction enters through one additive modulus. (a) \emph{American wedge}: on the undervalued legs $C^{\mathrm{am}}\ge C^{\mathrm{eu}}$ only helps, so only the two overvalued legs contribute, with $\varepsilon_\ell=\log(1+\bar\pi_\ell/C^{\mathrm{eu}}_\ell)\le\bar\pi_\ell/C^{\mathrm{eu}}_\ell$---recovering Corollary~\ref{cor:sharpcert} in log form. (b) \emph{Microstructure}: for mid-price $C$ and half-spread $s$, the ask and bid deviations are $\log(1+s/C)$ and $-\log(1-s/C)$ in absolute value; since the latter is larger, the conservative symmetric modulus is $\varepsilon_{\mathrm{spread}}=-\log(1-s/C)$ (directional one-sided terms may be kept separately, and the pipeline's ask/bid detection already applies them one-sidedly). (c) \emph{Strike rounding}: replacing the exact adjusted strike $\tilde K$ by a traded strike at distance $\Delta K$ perturbs the leg by $|\log C(\tilde K+\Delta K)-\log C(\tilde K)|\le\int_{\tilde K}^{\tilde K+\Delta K}|\partial_u\log C(u)|\,du\le|\Delta K|\sup_{u\in I}|\partial_u C(u)|/C(u)$, with $I$ the interval between the exact and rounded strike; the local form $\Delta K\,|\partial_K C|/C$ is a first-order approximation, not a bound. Summing (a)--(c) yields one additive margin covering early exercise, spreads, and discretization, with each term carrying the qualifications above.
\end{remark}

\section{The Direction of the American Perturbation}\label{sec:direction}

Write the American rectangle as European plus premia. The premium enters the \emph{overvalued} side at strike $\Kt1$ with weight $C(\Kt2,T_1)$ and the \emph{undervalued} side at strike $K_2$ with weight $C(K_1,T_1)$, plus terms in $\pi(\cdot,T_1)$ that vanish when no ex-date precedes $T_1$:
\[
\underbrace{\pi(\Kt1,T_2)\,C(\Kt2,T_1)}_{\text{pushes toward violation}}
\quad\text{vs.}\quad
\underbrace{\pi(K_2,T_2)\,C(K_1,T_1)}_{\text{pushes against}} .
\]

\begin{remark}[Observation: the American perturbation is not sign-determined by the dividend sandwich alone]\label{prop:direction}
Expanding the American price products introduces premium terms on both sides of the rectangle:
\[
\big(C_1^{\mathrm{eu}}+\pi_1\big)\big(C_2^{\mathrm{eu}}+\pi_2\big)\;-\;\big(C_3^{\mathrm{eu}}+\pi_3\big)\big(C_4^{\mathrm{eu}}+\pi_4\big),
\]
with legs $1{=}(K_1,T_1)$, $2{=}(K_2,T_2)$, $3{=}(\Kt1,T_2)$, $4{=}(\Kt2,T_1)$. The separate $\pi_i$ terms may enter with opposing signs, and without additional restrictions on the strike and maturity dependence of the early-exercise premium, the dividend sandwich $0\le\pi\le\bar D$ does not determine the sign of the net TP\textsubscript{2} perturbation. We do not claim monotonicity of $\pi$ in the strike.
\end{remark}

This observation is the precise reason detection alone cannot separate ``theory violation'' from ``American artifact'' in the uncertified regime---and hence the reason a learned, conditional treatment of dividend-window episodes (with ex-date timing, short-leg moneyness, and rate-regime features) is the natural methodology there.

\section{Toward Genuine American TP\textsubscript{2}: A Supermodular Optimal-Stopping Route}\label{sec:topkis}

The obstruction to extending Glasserman--Pirjol directly is that $\Cam(K,T)=\sup_{\tau\le T}\E[e^{-r\tau}(S_\tau-K)^+]$ is a \emph{value function}, not an expectation against a single terminal law. Log-supermodularity is not preserved by suprema in general---but it \emph{is} preserved by suprema of the right structure (Topkis's monotone value theorem: a supremum of a supermodular objective over an ascending constraint correspondence is supermodular). Since $\log \Cam(K,T)=\sup_{\tau\le T}\log\E[e^{-r\tau}(S_\tau-K)^+]$ (the logarithm being monotone), and writing $k=\log K$, TP\textsubscript{2} of $\Cam$ is supermodularity of the value function $(k,T)\mapsto\log \Cam(e^k,T)$, with the stopping-time lattice as the decision variable.

\begin{assumption}[Stopped tail dominance]\label{ass:stopped}
For $r=0$ and all bounded stopping times $\sigma\le\tau$ (a.s.), the pair $(S_\sigma,S_\tau)$ satisfies the GP conditional right-tail dominance condition on the relevant out-of-the-money region:
$\Prob(S_\sigma>y\mid S_\sigma>x)\;\le\;\Prob(S_\tau>y\mid S_\tau>x)$ for all $x_0\le x<y$.
\end{assumption}

\begin{conjecture}[Proposed route to American TP\textsubscript{2}]\label{thm:topkis}
Let $r=0$ and suppose Assumption~\ref{ass:stopped} holds, together with its lattice closure (the same condition for $\sigma\wedge\tau\le\sigma\vee\tau$ for arbitrary bounded stopping times $\sigma,\tau$). Then $(k,\tau)\mapsto\log\E[(S_\tau-e^{k})^+]$ has increasing differences in $(k,\tau)$ on the stopping-time lattice; the correspondence $T\mapsto\{\tau:\tau\le T\}$ is ascending; and by Topkis's theorem the value $\log\Cam(e^k,T)$ would be supermodular in $(k,T)$, i.e.\ American call prices would be TP\textsubscript{2} in $(K,T)$ on the corresponding out-of-the-money region.
\end{conjecture}

This is a proof program, not a theorem: it becomes one only when the following are established---a valid lattice structure for the admissible stopping rules; increasing differences jointly in the parameter and the stopping decision; closure under $\sigma\wedge\tau$ and $\sigma\vee\tau$; measurability and existence of optimizers; preservation of the stopped-tail condition; the positive-rate normalization; and the in-the-money mechanism. Assumption~\ref{ass:stopped} replaces GP's fixed-dates condition by its uniform version over stopping times. Optional sampling supplies martingale expectation identities under appropriate integrability conditions; it does not by itself establish the conditional tail-order properties required here. What makes this more than a reformulation is a reduction available in our setting:

\begin{proposition}[Discrete-dividend reduction]\label{prop:reduction}
For an American call on an asset paying proportional discrete dividends ($S_{t_i}=(1-q_i)S_{t_i^-}$) with $r=0$, optimal exercise is supported on $\{t_1^-,\ldots,t_n^-,T\}$, and the optimal rule at each $t_i^-$ is a threshold rule. Hence $\Cam(K,T)=\sup_{b\in\mathcal B}\,h(K,b;T)$ over a finite-dimensional family of threshold policies $b$ (the threshold vector is finite-dimensional; the set of admissible threshold values is generally uncountable), and Assumption~\ref{ass:stopped} needs to be verified only on the laws stopped by such policies. Moreover, each dividend map $S\mapsto(1-q_i)S$ \emph{preserves} the GP tail condition, which is invariant under multiplicative scaling: $\Prob(cS>y\mid cS>x)=\Prob(S>y/c\mid S>x/c)$.
\end{proposition}

\begin{corollary}[Conjectured; proof program]\label{cor:bsdiv}
In Black--Scholes with proportional discrete dividends and $r=0$, American call prices are TP\textsubscript{2} in $(K,T)$ on the out-of-the-money region: deterministic-date comparisons hold by Glasserman--Pirjol's Black--Scholes result, dividend maps preserve them by scale invariance, and the threshold-policy supremum preserves supermodularity by Topkis, \emph{provided} the increasing-differences verification across threshold policies (the lattice-closure step) goes through.
\end{corollary}

\noindent\textbf{What remains open.} (i) The lattice-closure verification in Proposition~\ref{prop:reduction}: increasing differences of $h$ in $(-k,b)$ across threshold policies---this is where a complete proof requires work beyond assembling known pieces; (ii) $r>0$, where the discount factor $e^{-r\tau}$ enters inside the expectation and the natural fix (a forward-measure normalization as in GP's carry adjustment) has to be carried through the stopping lattice; (iii) the in-the-money mechanism (GP Theorem~3.1(ii)) under stopping, which uses the increasing-convex order and does not compose with suprema as cleanly; (iv) cash-settled versus physically settled makes no difference to any statement here---only quoted prices enter---but \emph{negative} rates would reverse Lemma~\ref{lem:smalldiv}. We state Corollary~\ref{cor:bsdiv} as a conjecture with a proof program, not a theorem.

\section{Consequences for the Empirical Study}
\begin{enumerate}[leftmargin=1.4em,itemsep=1pt]
\item \textbf{Partition of episodes.} Every detected rectangle falls into one of three provable regimes: (R1) certified-European (Prop.~\ref{prop:window}: no ex-date in window, or \eqref{eq:smalldiv} holds for all four legs)---the SPX theory applies verbatim, and a violation there contradicts the GP conditions exactly as on SPX; (R2) dividend-robust (Prop.~\ref{prop:certificate}: violation margin exceeds $M$)---contradicts the European theory even granting the American premium its maximal effect; (R3) unresolved---the violation is within the early-exercise wedge, and Prop.~\ref{prop:direction} shows its sign is genuinely ambiguous. Only (R3) is theory-dark, and it is precisely where conditional machine-learned selection operates.
\item \textbf{Sharpened margins are computable per trade.} $M^{\mathrm{quote}}$ of Corollary~\ref{cor:sharpcert} needs only announced dividends and one quoted call at a lower strike (for \eqref{eq:doob}); it inherits the certificate's rigor, and the R2 (dividend-robust) region enlarges accordingly. $M^{\mathrm{model}}$ additionally uses the implied-volatility proxy (for \eqref{eq:chernoff}) and is diagnostic, not certifying.
\item \textbf{A rate-regime prediction.} By Remark~\ref{rem:regime}, the R1/R3 composition shifts with the interest-rate regime: at $r\approx0$ (2021) the sufficient condition fails for dividend-spanning contracts, while higher rates enlarge the certified set contract by contract. This is a falsifiable, theory-driven account of regime dependence that complements the label-base-rate shift documented in the paper, and motivates rate- and ex-date-timing features in the learned selectors.
\item \textbf{Feature-level grounding.} The margin $M$, the indicator of \eqref{eq:smalldiv}, the presence of an ex-date in $(0,T_1]$ versus $(T_1,T_2]$, and short-leg moneyness (which controls both assignment risk and the tightness of \eqref{eq:smalldiv}) are theory-mandated covariates of violation quality---a principled core for the feature dictionary, several members of which (ex-dividend timing flags, short-leg structure block) the empirical pipeline already contains.
\end{enumerate}

\begin{thebibliography}{9}
\bibitem{gp2026} P.~Glasserman and D.~Pirjol. When Are Option Prices TP\textsubscript{2}? \emph{Mathematical Finance} 36(3), 2026. \url{https://doi.org/10.1111/mafi.70024}
\bibitem{glp2025} P.~Glasserman, M.~Li, and D.~Pirjol. Trading TP\textsubscript{2} Option Violations. Working paper, SSRN 5392317, July 2025.
\bibitem{gp2023} P.~Glasserman and D.~Pirjol. Total Positivity and Relative Convexity of Option Prices. \emph{Frontiers of Mathematical Finance} 2(1):1--32, 2023.
\bibitem{merton} R.~C.~Merton. Theory of Rational Option Pricing. \emph{Bell Journal of Economics and Management Science} 4(1):141--183, 1973.
\bibitem{black1975} F.~Black. Fact and Fantasy in the Use of Options. \emph{Financial Analysts Journal} 31(4):36--41, 61--72, 1975.
\bibitem{vershynin} R.~Vershynin. \emph{High-Dimensional Probability: An Introduction with Applications in Data Science}. 2nd ed., Cambridge University Press, 2026.
\bibitem{topkis} D.~M.~Topkis. \emph{Supermodularity and Complementarity.} Princeton University Press, 1998.
\bibitem{cosma} A.~Cosma, S.~Galluccio, P.~Pederzoli, and O.~Scaillet. Early Exercise Decision in American Options with Dividends, Stochastic Volatility, and Jumps. \emph{JFQA} 55(1):331--356, 2020.
\end{thebibliography}


\appendix
\section{The Positive-Rate Doob Bound}\label{app:doobr}
Let $r>0$ and define the discounted dividend-reinvested process
\[
M_u\;=\;e^{-ru}S_u\;+\;\sum_{t_j\le u}e^{-rt_j}\delta_j,
\]
a $\Prob$-martingale under the risk-neutral measure. The target event satisfies
\[
\{S_{t_i^-}\ge K\}\;\subseteq\;\Big\{M_{t_i^-}\ge \lambda_i^*\Big\},
\qquad
\lambda_i^*\;=\;e^{-rt_i}K+\sum_{t_j<t_i}e^{-rt_j}\delta_j .
\]
For any $\lambda$ with $\sum_{t_j\le\tau}e^{-rt_j}\delta_j<\lambda<\lambda_i^*$ and any listed $\tau\ge t_i$, Doob's maximal inequality applied to the nonnegative submartingale $(M_u-\lambda)^+$ gives
\[
\Prob\big(S_{t_i^-}\ge K\big)\;\le\;\frac{\E\big[(M_\tau-\lambda)^+\big]}{\lambda_i^*-\lambda}.
\]
Finally,
\[
(M_\tau-\lambda)^+=e^{-r\tau}\Big(S_\tau-K'\Big)^+,
\qquad
K'\;=\;e^{r\tau}\lambda-\sum_{t_j\le\tau}e^{r(\tau-t_j)}\delta_j\;>\;0,
\]
so $\E[(M_\tau-\lambda)^+]=e^{-r\tau}\,\E[(S_\tau-K')^+]=C^{\mathrm{eu}}(K',\tau)\le C^{\mathrm{ask}}(K',\tau)$, where $C^{\mathrm{eu}}$ denotes the discounted (market-convention) European price. When $K'$ is not a listed strike, round it down to the nearest listed strike; the call price is nonincreasing in strike, so the bound remains valid. Optimizing over admissible $\lambda$ and listed $\tau$ gives the positive-rate analogue of \eqref{eq:doob}.

\end{document}
